\documentclass[10pt,letterpaper]{article}
\usepackage{cornell-notes}
\usepackage{mathtools}

\title{Chapter 8: Sorting and Selection}
\author{}
\date{}

\setDocTitle{Chapter 8: Sorting and Selection}
\setDocSubtitle{Numerical Methods}
\setDocVersion{v1.0}
\setDocStatus{Study Companion}
\setDocOwner{Numerical Methods Cornell Notes}
\setDocAuthor{}
\setDocDate{\today}
\setDocSummary{Cornell-format study and review notes for Chapter 8: Sorting and Selection.}

\setCornellCollection{Numerical Methods Cornell Notes}
\setCornellUnitType{Chapter}
\setCornellUnitNumber{8}
\setCornellUnitTitle{Sorting and Selection}
\setCornellCompanionLabel{Study and review companion}

\hypersetup{pdftitle={Chapter 8: Sorting and Selection},pdfsubject={Cornell notes},pdfauthor={}}

\begin{document}
\maketitle

\section*{Chapter Overview}
This chapter organizes the mathematical and computational ideas behind sorting and selection. The notes preserve the numbered section sequence while emphasizing method selection, explicit assumptions, numerical reliability, cost, and verification.

\section*{Learning Objectives}
Explain the governing problem class; compare Straight Insertion and Shell's Method, Quicksort, Heapsort, and Indexing and Ranking; design defensible stopping and verification checks; distinguish problem conditioning from algorithmic stability.

\section*{Prerequisites}
Arrays, recursion, invariants, comparison functions, and asymptotic complexity.

\section*{Operational Workflow}
Define key ordering and stability needs; inspect input size and presortedness; select sort, indirect ordering, selection, or disjoint-set processing; validate permutation and boundary invariants.

\subsection*{Formula and notation anchor}
\[
T_{\mathrm{quick}}(n)=T(k)+T(n-k-1)+\Theta(n)
\]
Balanced partitions give $\Theta(n\log n)$ work; repeatedly extreme partitions give $\Theta(n^2)$.

\begin{warnbox}{Numerical warning}
Inconsistent comparators, worst-case pivots, deep recursion, unstable reordering, integer overflow in index arithmetic, and destructive partitioning of caller data.
\end{warnbox}

\section{8.0 Introduction}
\begin{cornell}
\noterow{What is the central idea?}{Ordering and selection algorithms trade comparisons, swaps, auxiliary storage, stability, adaptivity, and worst-case guarantees.}
\noterow{How should it be used?}{For Introduction, begin from explicit inputs, outputs, preconditions, and representation choices.  Define key ordering and stability needs; inspect input size and presortedness; select sort, indirect ordering, selection, or disjoint-set processing; validate permutation and boundary invariants.}
\noterow{Cost and reliability}{Analyze arithmetic work, storage, data movement, and convergence separately.  Relevant failure pressures include inconsistent comparators, worst-case pivots, deep recursion, unstable reordering, integer overflow in index arithmetic, and destructive partitioning of caller data.  Use scaled tolerances and report both success criteria and diagnostic evidence.}
\noterow{Implementation checklist}{\begin{itemize}
\item Validate dimensions, domains, ordering, and structural assumptions.
\item Guard small divisors, invalid states, overflow, underflow, and nonfinite values.
\item Make mutation, workspace, indexing, and termination behavior explicit.
\item Test a simple exact case, a difficult valid case, a boundary case, and an expected failure.
\end{itemize}}
\end{cornell}
\begin{summarybox}{Section summary}
Ordering and selection algorithms trade comparisons, swaps, auxiliary storage, stability, adaptivity, and worst-case guarantees.  Select this technique only when its assumptions match the data, and accept a result only after an independent residual, error estimate, invariant, or refinement check.
\end{summarybox}

\section{8.1 Straight Insertion and Shell's Method}
\begin{cornell}
\noterow{What is the central idea?}{Insertion sort excels on small or nearly ordered inputs; Shell sorting reduces long-distance disorder through a decreasing gap schedule.}
\noterow{How should it be used?}{For Straight Insertion and Shell's Method, begin from explicit inputs, outputs, preconditions, and representation choices.  Define key ordering and stability needs; inspect input size and presortedness; select sort, indirect ordering, selection, or disjoint-set processing; validate permutation and boundary invariants.}
\noterow{Quantitative lens}{Identify the governing equation, recurrence, probability law, or geometric predicate for Straight Insertion and Shell's Method.  Define every variable and scale; then derive a residual, error estimate, or invariant that can be checked independently.}
\noterow{Cost and reliability}{Analyze arithmetic work, storage, data movement, and convergence separately.  Relevant failure pressures include inconsistent comparators, worst-case pivots, deep recursion, unstable reordering, integer overflow in index arithmetic, and destructive partitioning of caller data.  Use scaled tolerances and report both success criteria and diagnostic evidence.}
\noterow{Implementation checklist}{\begin{itemize}
\item Validate dimensions, domains, ordering, and structural assumptions.
\item Guard small divisors, invalid states, overflow, underflow, and nonfinite values.
\item Make mutation, workspace, indexing, and termination behavior explicit.
\item Test a simple exact case, a difficult valid case, a boundary case, and an expected failure.
\end{itemize}}
\end{cornell}
\begin{summarybox}{Section summary}
Insertion sort excels on small or nearly ordered inputs; Shell sorting reduces long-distance disorder through a decreasing gap schedule.  Select this technique only when its assumptions match the data, and accept a result only after an independent residual, error estimate, invariant, or refinement check.
\end{summarybox}

\section{8.2 Quicksort}
\begin{cornell}
\noterow{What is the central idea?}{Quicksort partitions around a pivot and recurses on subarrays, delivering excellent average performance with pivot and recursion safeguards.}
\noterow{How should it be used?}{For Quicksort, begin from explicit inputs, outputs, preconditions, and representation choices.  Define key ordering and stability needs; inspect input size and presortedness; select sort, indirect ordering, selection, or disjoint-set processing; validate permutation and boundary invariants.}
\noterow{Quantitative anchor}{\[
A=A_{<p}\,\Vert\,A_{=p}\,\Vert\,A_{>p}
\]
State the dimensions, scaling, and validity assumptions before using this relation.  Check the resulting residual or invariant rather than trusting an iteration count alone.}
\noterow{Cost and reliability}{Analyze arithmetic work, storage, data movement, and convergence separately.  Relevant failure pressures include inconsistent comparators, worst-case pivots, deep recursion, unstable reordering, integer overflow in index arithmetic, and destructive partitioning of caller data.  Use scaled tolerances and report both success criteria and diagnostic evidence.}
\noterow{Implementation checklist}{\begin{itemize}
\item Validate dimensions, domains, ordering, and structural assumptions.
\item Guard small divisors, invalid states, overflow, underflow, and nonfinite values.
\item Make mutation, workspace, indexing, and termination behavior explicit.
\item Test a simple exact case, a difficult valid case, a boundary case, and an expected failure.
\end{itemize}}
\end{cornell}
\begin{summarybox}{Section summary}
Quicksort partitions around a pivot and recurses on subarrays, delivering excellent average performance with pivot and recursion safeguards.  Select this technique only when its assumptions match the data, and accept a result only after an independent residual, error estimate, invariant, or refinement check.
\end{summarybox}

\section{8.3 Heapsort}
\begin{cornell}
\noterow{What is the central idea?}{Heapsort maintains a heap invariant to select successive extremes, guaranteeing in-place worst-case logarithmic work per extraction.}
\noterow{How should it be used?}{For Heapsort, begin from explicit inputs, outputs, preconditions, and representation choices.  Define key ordering and stability needs; inspect input size and presortedness; select sort, indirect ordering, selection, or disjoint-set processing; validate permutation and boundary invariants.}
\noterow{Quantitative lens}{Identify the governing equation, recurrence, probability law, or geometric predicate for Heapsort.  Define every variable and scale; then derive a residual, error estimate, or invariant that can be checked independently.}
\noterow{Cost and reliability}{Analyze arithmetic work, storage, data movement, and convergence separately.  Relevant failure pressures include inconsistent comparators, worst-case pivots, deep recursion, unstable reordering, integer overflow in index arithmetic, and destructive partitioning of caller data.  Use scaled tolerances and report both success criteria and diagnostic evidence.}
\noterow{Implementation checklist}{\begin{itemize}
\item Validate dimensions, domains, ordering, and structural assumptions.
\item Guard small divisors, invalid states, overflow, underflow, and nonfinite values.
\item Make mutation, workspace, indexing, and termination behavior explicit.
\item Test a simple exact case, a difficult valid case, a boundary case, and an expected failure.
\end{itemize}}
\end{cornell}
\begin{summarybox}{Section summary}
Heapsort maintains a heap invariant to select successive extremes, guaranteeing in-place worst-case logarithmic work per extraction.  Select this technique only when its assumptions match the data, and accept a result only after an independent residual, error estimate, invariant, or refinement check.
\end{summarybox}

\section{8.4 Indexing and Ranking}
\begin{cornell}
\noterow{What is the central idea?}{An index permutation records sorted order without moving large records; ranks provide the inverse mapping from items to order positions.}
\noterow{How should it be used?}{For Indexing and Ranking, begin from explicit inputs, outputs, preconditions, and representation choices.  Define key ordering and stability needs; inspect input size and presortedness; select sort, indirect ordering, selection, or disjoint-set processing; validate permutation and boundary invariants.}
\noterow{Quantitative lens}{Identify the governing equation, recurrence, probability law, or geometric predicate for Indexing and Ranking.  Define every variable and scale; then derive a residual, error estimate, or invariant that can be checked independently.}
\noterow{Cost and reliability}{Analyze arithmetic work, storage, data movement, and convergence separately.  Relevant failure pressures include inconsistent comparators, worst-case pivots, deep recursion, unstable reordering, integer overflow in index arithmetic, and destructive partitioning of caller data.  Use scaled tolerances and report both success criteria and diagnostic evidence.}
\noterow{Implementation checklist}{\begin{itemize}
\item Validate dimensions, domains, ordering, and structural assumptions.
\item Guard small divisors, invalid states, overflow, underflow, and nonfinite values.
\item Make mutation, workspace, indexing, and termination behavior explicit.
\item Test a simple exact case, a difficult valid case, a boundary case, and an expected failure.
\end{itemize}}
\end{cornell}
\begin{summarybox}{Section summary}
An index permutation records sorted order without moving large records; ranks provide the inverse mapping from items to order positions.  Select this technique only when its assumptions match the data, and accept a result only after an independent residual, error estimate, invariant, or refinement check.
\end{summarybox}

\section{8.5 Selecting the Mth Largest}
\begin{cornell}
\noterow{What is the central idea?}{Partition-based selection discards the side that cannot contain the target rank and usually avoids fully sorting the data.}
\noterow{How should it be used?}{For Selecting the Mth Largest, begin from explicit inputs, outputs, preconditions, and representation choices.  Define key ordering and stability needs; inspect input size and presortedness; select sort, indirect ordering, selection, or disjoint-set processing; validate permutation and boundary invariants.}
\noterow{Quantitative lens}{Identify the governing equation, recurrence, probability law, or geometric predicate for Selecting the Mth Largest.  Define every variable and scale; then derive a residual, error estimate, or invariant that can be checked independently.}
\noterow{Cost and reliability}{Analyze arithmetic work, storage, data movement, and convergence separately.  Relevant failure pressures include inconsistent comparators, worst-case pivots, deep recursion, unstable reordering, integer overflow in index arithmetic, and destructive partitioning of caller data.  Use scaled tolerances and report both success criteria and diagnostic evidence.}
\noterow{Implementation checklist}{\begin{itemize}
\item Validate dimensions, domains, ordering, and structural assumptions.
\item Guard small divisors, invalid states, overflow, underflow, and nonfinite values.
\item Make mutation, workspace, indexing, and termination behavior explicit.
\item Test a simple exact case, a difficult valid case, a boundary case, and an expected failure.
\end{itemize}}
\end{cornell}
\begin{summarybox}{Section summary}
Partition-based selection discards the side that cannot contain the target rank and usually avoids fully sorting the data.  Select this technique only when its assumptions match the data, and accept a result only after an independent residual, error estimate, invariant, or refinement check.
\end{summarybox}

\section{8.6 Determination of Equivalence Classes}
\begin{cornell}
\noterow{What is the central idea?}{Union-find maintains disjoint sets using path compression and rank or size heuristics, making repeated merges and queries nearly constant amortized time.}
\noterow{How should it be used?}{For Determination of Equivalence Classes, begin from explicit inputs, outputs, preconditions, and representation choices.  Define key ordering and stability needs; inspect input size and presortedness; select sort, indirect ordering, selection, or disjoint-set processing; validate permutation and boundary invariants.}
\noterow{Quantitative anchor}{\[
\operatorname{find}(x)\to\text{representative},\qquad \operatorname{union}(a,b)
\]
State the dimensions, scaling, and validity assumptions before using this relation.  Check the resulting residual or invariant rather than trusting an iteration count alone.}
\noterow{Cost and reliability}{Analyze arithmetic work, storage, data movement, and convergence separately.  Relevant failure pressures include inconsistent comparators, worst-case pivots, deep recursion, unstable reordering, integer overflow in index arithmetic, and destructive partitioning of caller data.  Use scaled tolerances and report both success criteria and diagnostic evidence.}
\noterow{Implementation checklist}{\begin{itemize}
\item Validate dimensions, domains, ordering, and structural assumptions.
\item Guard small divisors, invalid states, overflow, underflow, and nonfinite values.
\item Make mutation, workspace, indexing, and termination behavior explicit.
\item Test a simple exact case, a difficult valid case, a boundary case, and an expected failure.
\end{itemize}}
\end{cornell}
\begin{summarybox}{Section summary}
Union-find maintains disjoint sets using path compression and rank or size heuristics, making repeated merges and queries nearly constant amortized time.  Select this technique only when its assumptions match the data, and accept a result only after an independent residual, error estimate, invariant, or refinement check.
\end{summarybox}

\section*{Method comparison}
\begin{center}
\small
\begin{tabularx}{\linewidth}{>{\bfseries\raggedright\arraybackslash}p{0.22\linewidth}XX}
\toprule
Method or lens & Best use & Main caution \\
\midrule
Insertion sort & Small or nearly sorted arrays & Stable and simple; quadratic worst case \\
Quicksort & General in-memory arrays & Fast average; guard pivot and recursion \\
Heapsort & In-place worst-case guarantee & Predictable $n\log n$; weaker locality \\
\bottomrule
\end{tabularx}
\end{center}

\section*{Worked example}
\begin{exambox}{Independent worked example}
To select the third-smallest value in $[8,2,5,1,7]$, partition around $5$ to obtain values below it, the pivot, and values above it.  Two values lie below, so $5$ has target rank and no full sort is needed.  Duplicate keys require a three-way partition to prevent ambiguous progress.
\end{exambox}

\section*{Chapter synthesis}
\begin{summarybox}{Chapter synthesis}
The chapter's methods should be viewed as a decision system rather than a list of routines.  Begin with the mathematical problem and its conditioning, exploit structure, select an algorithm with compatible stability and cost, and verify the computed answer with evidence that is independent of the internal iteration.  The recurring implementation discipline is: define key ordering and stability needs; inspect input size and presortedness; select sort, indirect ordering, selection, or disjoint-set processing; validate permutation and boundary invariants.
\end{summarybox}

\subsection*{Common mistakes and numerical hazards}
\begin{warnbox}{Common mistakes and numerical hazards}
Inconsistent comparators, worst-case pivots, deep recursion, unstable reordering, integer overflow in index arithmetic, and destructive partitioning of caller data.  A second recurring mistake is reporting only a computed value: always retain tolerances, iteration counts, residuals, refinement behavior, and relevant structural checks.
\end{warnbox}

\subsection*{Retrieval practice}
\textbf{Questions}
\begin{enumerate}
\item What problem does Straight Insertion and Shell's Method address, and which assumption is easiest to violate?
\item What problem does Quicksort address, and which assumption is easiest to violate?
\item What problem does Heapsort address, and which assumption is easiest to violate?
\item What problem does Selecting the Mth Largest address, and which assumption is easiest to violate?
\item What problem does Determination of Equivalence Classes address, and which assumption is easiest to violate?
\item How do conditioning, stability, and the final residual play different roles in this chapter?
\end{enumerate}

\textbf{Brief answers}
\begin{enumerate}
\item Insertion sort excels on small or nearly ordered inputs; Shell sorting reduces long-distance disorder through a decreasing gap schedule. Recheck the section's domain, structure, scaling, and failure diagnostics.
\item Quicksort partitions around a pivot and recurses on subarrays, delivering excellent average performance with pivot and recursion safeguards. Recheck the section's domain, structure, scaling, and failure diagnostics.
\item Heapsort maintains a heap invariant to select successive extremes, guaranteeing in-place worst-case logarithmic work per extraction. Recheck the section's domain, structure, scaling, and failure diagnostics.
\item Partition-based selection discards the side that cannot contain the target rank and usually avoids fully sorting the data. Recheck the section's domain, structure, scaling, and failure diagnostics.
\item Union-find maintains disjoint sets using path compression and rank or size heuristics, making repeated merges and queries nearly constant amortized time. Recheck the section's domain, structure, scaling, and failure diagnostics.
\item Conditioning describes sensitivity of the mathematical problem, stability describes error propagation by the algorithm, and the residual measures satisfaction of the computed equations; none is a universal substitute for the others.
\end{enumerate}

\subsection*{Mastery checklist}
\begin{itemize}
\item[$\square$] I can explain and select Straight Insertion and Shell's Method without relying on a memorized code listing.
\item[$\square$] I can explain and select Quicksort without relying on a memorized code listing.
\item[$\square$] I can explain and select Heapsort without relying on a memorized code listing.
\item[$\square$] I can explain and select Indexing and Ranking without relying on a memorized code listing.
\item[$\square$] I can explain and select Selecting the Mth Largest without relying on a memorized code listing.
\item[$\square$] I can state a scaled stopping rule and an independent verification check.
\item[$\square$] I can identify at least one conditioning risk and one algorithmic stability risk.
\item[$\square$] I can estimate time, storage, and data-movement costs at the appropriate level.
\end{itemize}

\end{document}

